\documentclass[11pt,a4paper]{article}
\usepackage[T1]{fontenc}
\usepackage{mathptmx}
\usepackage[scaled=0.92]{helvet}
\usepackage{courier}
\usepackage[utf8]{inputenc}
\usepackage[brazil]{babel}
\usepackage[margin=2.2cm]{geometry}
\usepackage{booktabs,graphicx,caption,float,amsmath,amssymb,bm,xcolor,array,enumitem,tikz,longtable}
\usepackage[hidelinks]{hyperref}
\usetikzlibrary{arrows.meta,positioning,calc}
\setlist{itemsep=1pt, topsep=3pt}
\setlength{\emergencystretch}{3em}
\sloppy
\makeatletter\g@addto@macro\UrlBreaks{\do\_}\makeatother
\newcolumntype{L}[1]{>{\raggedright\arraybackslash}p{#1}}

\title{Guia de estudo\\[3pt]\large Talude não saturado sob chuva com campos aleatórios no NeoPZ:\\
  formulação, verificação, estabilidade, incerteza e confiabilidade}
\author{}
\date{14 de novembro de 2026}

\begin{document}
\maketitle
\tableofcontents

\section{O mapa do trabalho}
\label{sec:mapa}

\paragraph{O objetivo.} Calcular a probabilidade de ruptura de um talude não saturado, sob chuva ou rebaixamento, com a
resistência e a permeabilidade como campos aleatórios. Os alvos são reproduzir a Seção 6 de Vargas Ceron et al.\ (2025)
e, depois, a chuva de Rahardjo et al.\ (2007), com o modelo acoplado do NeoPZ.

\paragraph{A cadeia.} Cada elo depende do anterior e tem um documento e um código.
\begin{itemize}
  \item a física do solo não saturado;
  \item a discretização e a sua verificação;
  \item o fator de segurança;
  \item o cenário hidráulico;
  \item a incerteza espacial;
  \item a propagação por Monte Carlo;
  \item a probabilidade de falha.
\end{itemize}

\begin{figure}[H]
\centering
\begin{tikzpicture}[
  caixa/.style={draw, rounded corners=2pt, fill=#1, text width=15.2cm, inner sep=4pt, font=\small, align=left},
  caixa/.default={gray!8}, seta/.style={-Stealth, thick, gray!70!black}, node distance=2.5mm]
  \node[caixa=blue!6] (f) {\textbf{1.\ Formulação $\bm u$--$p_w$--$p_a$}: Bishop, retenção, Darcy, Mohr--Coulomb PV2,
    Euler implícito, Newton.\\
    \emph{Documentos:} \texttt{formulacao}, \path{formulacao_verick}, \path{resumo_nao_saturado}.
    \emph{Código:} \texttt{TPZMatUnsaturatedUPwPa}, \texttt{TPZUnsatAnalysis} (Seção~\ref{sec:formulacao}).};
  \node[caixa=blue!6, below=of f] (v) {\textbf{2.\ Verificação}: ponto material, água, ar, Terzaghi, Mandel,
    Liakopoulos, talude seco contra Bishop.\\
    \emph{Documentos:} revisão bibliográfica e resumo com Liakopoulos. \emph{Código:} \texttt{Unsat*},
    \texttt{LearningLiakopoulos} (Seção~\ref{sec:verificacao}).};
  \node[caixa=green!6, below=of v] (s) {\textbf{3.\ Fator de segurança}: SRM e GIM com as pressões congeladas;
    arc-length.\\
    \emph{Documento:} \path{relatorio_fs}. \emph{Código:} \texttt{TPZUnsatAnalysis::ArcLength*},
    \path{Common/SlopeSolver} (Seção~\ref{sec:fs}).};
  \node[caixa=green!6, below=of s] (c) {\textbf{4.\ Cenário hidráulico}: chuva (Rahardjo, transiente de Richards) ou
    rebaixamento (Ceron, percolação permanente).\\
    \emph{Documentos:} \path{relatorio_chuva} e o artigo da intensidade crítica. \emph{Código:}
    \path{Common/SlopeRain}, \texttt{SlopeModel}, \texttt{SlopeRandomField} (Seção~\ref{sec:chuva}).};
  \node[caixa=orange!8, below=of c] (k) {\textbf{5.\ Campos aleatórios}: Karhunen--Loève de $c'$, $\phi'$ e $k$:
    forma fechada, Galerkin e Nyström.\\
    \emph{Documentos:} o relatório de 11/11 (\texttt{relatorio-2}) e o \texttt{nymstron.nb}. \emph{Código:}
    \texttt{KLSeparavel.h}, \path{Common/SlopeFields}, \texttt{KLCilindroMika}, \texttt{VerifyRandomFields}
    (Seção~\ref{sec:campos}).};
  \node[caixa=orange!8, below=of k] (m) {\textbf{6.\ Monte Carlo}: refinamento por realização, níveis de malha,
    campanhas no fermi.\\
    \emph{Documento:} o relatório de 11/11. \emph{Código:} \texttt{SlopeRandomFieldRefina}, \texttt{SlopeMonteCarlo}
    (Seção~\ref{sec:mc}).};
  \node[caixa=red!6, below=of m] (r) {\textbf{7.\ Confiabilidade}: $P_f$ e o seu CoV; FORM e amostragem por
    importância para $P_f$ pequenos.\\
    \emph{Documento:} \texttt{confiabilidade} (pasta \path{form_MIS}). \emph{Código:} \texttt{Confiabilidade}
    (Seção~\ref{sec:confiab}).};
  \foreach \a/\b in {f/v, v/s, s/c, c/k, k/m, m/r} \draw[seta] (\a.south) -- (\b.north);
\end{tikzpicture}
\caption{A cadeia do trabalho. Azul: o modelo; verde: a estabilidade sob um cenário; laranja: a incerteza;
  vermelho: o resultado.}
\label{fig:mapa}
\end{figure}

\paragraph{Os documentos do zip e onde estão no repositório.}
\begin{center}\small
\begin{longtable}{@{}L{5.6cm} L{2.6cm} L{7.4cm}@{}}
\toprule
Arquivo do zip & Data, páginas & No repositório \\
\midrule
\endhead
\path{formulacao/formulacao.tex} & 30/10, 25 p. & \path{docs/biot_nao_saturado/formulacao} (idêntico) \\
\path{formulacao/formulacao_verick.tex} & 01/11, 31 p. & não está \\
\path{formulacao/resumo_nao_saturado.tex} & 01/11, 13 p. & não está (versão curta) \\
\path{Liakopoulos/resumo/resumo_nao_saturado.tex} & 02/11, 24 p. & \path{docs/biot_nao_saturado/resumo} (idêntico) \\
\path{revisao_bibliografica_nao_saturado_neopz.pdf} & 31/10, 20 p. & não está \\
\path{arclength/relatorio_fs.tex} & 07/11, 34 p. & \path{docs/fs_arc_length} (idêntico) \\
\path{relatorio_form/relatorio_chuva.tex} & 14/11, 12 p. & não está, nem o código dele \\
\path{formulacao_random_fields/relatorio-2.pdf} & 11/11, 36 p. & \path{SlopeRandomFieldRefina/relatorio} \\
\path{formulacao_random_fields/nymstron.nb} & 14/11 & não está (o seu caderno do Nyström) \\
\path{form_MIS/confiabilidade.*} & 14/11, 19 p. & \path{Projects/Confiabilidade}. Na sua cópia do \texttt{.wl}, o esboço
  \texttt{gTalude} da Seção 7 está descomentado: só define a função, e chamá-la exige um executável que ainda não existe \\
\path{depricated/} & 16/10 a 30/10 & a evolução da formulação (Apêndice~\ref{sec:historia}) \\
\bottomrule
\end{longtable}
\end{center}

\section{O problema-alvo e os números de referência}
\label{sec:alvo}

Dois problemas de referência guiam o trabalho. O primeiro dá a parte estocástica, o segundo a parte da chuva.

\paragraph{Vargas Ceron et al.\ (2025), Seção 6: o talude com campos aleatórios.}
\begin{itemize}
  \item \textbf{Geometria e carregamento:} talude de $H = 5$\,m e $45^\circ$, com rebaixamento total ($h_w = H$).
  \item \textbf{Cálculo do fator:} percolação permanente com o $k_v$ da realização, e $\Gamma$ pelo método de aumento
    da gravidade (GIM), com as pressões congeladas.
  \item \textbf{Campos lognormais independentes} (Tabela 2):
    \begin{itemize}
      \item $c'$: 10\,kPa, CoV 30\,\%;
      \item $\phi'$: $30^\circ$, CoV 10\,\%;
      \item $k_v$: CoV 60\,\%.
    \end{itemize}
    Os três têm o núcleo exponencial separável, com $L_x = 20$\,m e $L_y = 2$\,m, e são gerados por Karhunen-Loève com
    $M = 2000$ termos.
  \item \textbf{Ruptura e probabilidade:} a ruptura é $\Gamma < 1$, e $P_f = n_f/S$.
  \item \textbf{Resultados do artigo:} com $S = 10\,000$, $P_f = 11{,}50\,\%$ (CoV 2,77\,\%). $\Gamma$ é lognormal, com
    média 1,353, mediana 1,313 e CoV 23,5\,\%.
\end{itemize}

\paragraph{Rahardjo et al.\ (2007), Série A: o talude sob chuva.}
\begin{itemize}
  \item \textbf{Geometria e solo:} talude de $45^\circ$, $c' = 10$\,kPa, $\phi' = \phi^b = 26^\circ$, $\gamma = 20$\,kN/m$^3$
    e $n = 0{,}45$.
  \item \textbf{Hidráulica:} curva de retenção de Fredlund--Xing com $a = 10$, 50 ou 100\,kPa e $k_s = 10^{-4}$, $10^{-5}$
    ou $10^{-6}$\,m/s.
  \item \textbf{Uso:} é o transiente de chuva do \texttt{SlopeRandomField}, do relatório da chuva
    (Seção~\ref{sec:chuva}) e do artigo em preparação sobre a intensidade crítica de chuva.
\end{itemize}

\paragraph{Onde estamos} (números do repositório, branch \texttt{UnsaturatedSoil}).
\begin{itemize}
  \item \textbf{Determinístico do Ceron} (\texttt{SlopeMonteCarlo}, GIM adaptativo hp):
    \begin{itemize}
      \item nos níveis 0 a 4, $\Gamma = 1{,}7444$; 1,5559; 1,4488; 1,4032; 1,3850, com 126 a 1314 triângulos;
      \item a sequência extrapola para cerca de 1,373, o mesmo $\Gamma$ que o método do Ceron dá neste domínio (1,374);
      \item o artigo dá 1,336. Segundo o README do \texttt{SlopeMonteCarlo}, a diferença vem do contorno distante da
        percolação e não muda com a malha.
    \end{itemize}
  \item \textbf{Monte Carlo:}
    \begin{itemize}
      \item com 3 níveis e 1000 realizações, $P_f = 10{,}3\,\%$ (o artigo dá 11,5\,\%), com média lognormal de 1,398;
      \item no estudo de refinamento do \texttt{SlopeRandomFieldRefina}, $P_f = 10{,}8\,\%$ na malha de referência
        \texttt{foco125}, com CoV de 9,1\,\% para 1000 realizações;
      \item a campanha de 10\,000 realizações com 4 níveis, cerca de 95\,s por tarefa no fermi, está pronta para rodar.
    \end{itemize}
\end{itemize}

\section{Formulação do solo não saturado (modelo \texorpdfstring{$\bm u$--$p_w$--$p_a$}{u-pw-pa})}
\label{sec:formulacao}

\paragraph{Os documentos.}
\begin{itemize}
  \item \textbf{\texttt{formulacao.tex}} (30/10, 25 páginas; igual a \path{docs/biot_nao_saturado/formulacao}): o
    documento de trabalho das Fases 1 e 2. Para cada bloco ele traz as equações, a discretização, os resíduos e o
    Jacobiano, cita o arquivo e as linhas do código e numera as discrepâncias (D1 a D20).
  \item \textbf{\path{formulacao_verick.tex}} (01/11, 31 páginas; só no zip): o mesmo conteúdo, na ordem e na notação
    da tese de Santos (2009), \emph{Simulador de meios porosos saturados elastoplásticos}. Acrescenta a contagem das
    25 variáveis e 25 equações, o problema poroelastoplástico explícito e a D21 (desacoplamento). Os números de
    verificação são os mesmos.
  \item \textbf{\path{resumo_nao_saturado.tex}}: a versão compacta (01/11, 13 páginas, com a tabela de
    correspondência de notação) e a versão com o exemplo de Liakopoulos e o projeto \texttt{LearningLiakopoulos}
    (02/11, 24 páginas; igual a \path{docs/biot_nao_saturado/resumo}).
  \item \textbf{Fontes:} as notas [N] (\path{biot_nao_saturado_agua_ar3} e \texttt{ar4}) e o artigo [A], Lira
    Cecílio (2026), \emph{J.\ Eng.\ Math.}\ 157:10, com o retorno analítico do Mohr--Coulomb PV2.
\end{itemize}
Cuidado com a notação, que muda entre os documentos (tese de Santos $\leftrightarrow$ \texttt{formulacao.tex}):
porosidade $\phi \leftrightarrow n$, atrito $\vartheta \leftrightarrow \phi'$, $C_s \leftrightarrow \Gamma$,
$1/N \leftrightarrow \eta$, $K_a \leftrightarrow P$ e $\Delta\zeta_{sk} \leftrightarrow \Theta$. Abaixo uso a do
\texttt{formulacao.tex}, que é a do cabeçalho do material.

\paragraph{Hipóteses.}
\begin{itemize}
  \item \textbf{Regime:} pequenas deformações, quase estático, isotérmico; água e ar imiscíveis, sem mudança de fase;
    Darcy com permeabilidades relativas; estado plano de deformação.
  \item \textbf{Sinais:} pressões manométricas positivas em compressão; tensões positivas em tração.
  \item \textbf{Constantes:} $\alpha$, $K_s$, $n$, $\bm K$ e $\gamma$ são constantes, e não há rigidez geométrica.
  \item \textbf{Ar:} $p_a$ é incógnita em todo o domínio.
\end{itemize}

\paragraph{Leis constitutivas.}
\begin{itemize}
  \item \textbf{Tensão efetiva de Bishop:} $\bm\sigma = \bm\sigma' - \alpha p_B \bm I$, com
    $p_B = \chi p_w + (1 - \chi) p_a$ e $p_c = p_a - p_w$.
    \begin{itemize}
      \item O padrão do material é $\chi = S_w$, que dá $p_B = p_a - S_w p_c$.
      \item As aplicações de talude usam $\chi$ constante, igual a 1 ou a $\tan\phi^b/\tan\phi'$.
    \end{itemize}
  \item \textbf{Retenção} $S_w(p_c)$, com $\Gamma = -dS_w/dp_c \geq 0$: van Genuchten, Fredlund--Xing (o padrão do
    código, com $a = 50$\,kPa), Liakopoulos e Gardner.
  \item \textbf{Permeabilidades relativas:} Mualem, $S_e^p$ e Brooks--Corey, com $k_{rw} \geq 10^{-12}$ e
    $k_{ra} \geq 10^{-4}$.
  \item \textbf{Fluidos:} a água com módulo $K_w$; o ar como gás ideal isotérmico, com $K_a = P = p_a + p_{atm}$.
  \item \textbf{Esqueleto:} Mohr--Coulomb associado ($\psi = \phi'$), o PV2 do artigo [A].
    \begin{itemize}
      \item Retorno analítico por setor: elástico, face, aresta direita, aresta esquerda e ápice.
      \item Tangente consistente em forma fechada.
      \item Com $\psi \neq \phi'$, o código para (\texttt{DebugStop}).
    \end{itemize}
\end{itemize}

\paragraph{Equações de campo.}
\begin{gather}
  \operatorname{div}\bm\sigma + \rho_b \bm g = \bm 0, \qquad
  \bm q_\pi = -\lambda_\pi \bm K (\nabla p_\pi - \gamma_\pi \bm g), \quad
  \lambda_w = \frac{k_{rw}}{\gamma_w}, \ \lambda_a = \frac{\mu_w}{\mu_a}\frac{k_{ra}}{\gamma_w}, \\
  \dot\zeta_\pi = S_\pi(\alpha\dot\varepsilon_v + \eta\dot p_B) + n\dot S_\pi + \frac{nS_\pi}{K_\pi}\dot p_\pi, \qquad
  \dot\zeta_\pi + \operatorname{div}\bm q_\pi + c_\pi\, \bm q_\pi\cdot\nabla p_\pi = 0, \qquad \pi \in \{w, a\},
\end{gather}
com $\eta = (\alpha - n)/K_s$ e $c_a = 1/P$. O termo de $c_w$ não existe no código. Os coeficientes tangentes de
Bishop são $\chi_w = \partial p_B/\partial p_w = S_w - \Gamma p_c$ e $\chi_a = S_a + \Gamma p_c$, com soma 1.

\paragraph{Discretização.}
\begin{itemize}
  \item \textbf{Espaço:} Taylor--Hood: $\bm u$ em P2 e $p_w$, $p_a$ em P1, que satisfaz a condição inf-sup. A regra tem
    3 pontos.
  \item \textbf{Armazenamento:} as parcelas de saturação e de compressibilidade do ar ficam concentradas nos vértices,
    com $m_j = \int N_j$.
  \item \textbf{Tempo:} Euler implícito, monolítico. As equações de massa são multiplicadas por $\Delta t$, então
    $\Delta t = 0$ vale e é usado no estado geostático e no FS, com as pressões congeladas.
  \item \textbf{Armazenamento secante}, que é conservativo:
    \begin{equation}
      \Delta\zeta_w = S_w^{n+1}\Theta + n\Delta S_w + \frac{nS_w^{n+1}}{K_w}\Delta p_w, \qquad
      \Delta\zeta_a = S_a^{n+1}\Theta - n\Delta S_w + \frac{nS_a^{n}}{P^{n+1}}\Delta p_a, \qquad
      \Theta = \alpha\Delta\varepsilon_v + \eta\Delta p_B .
    \end{equation}
    \begin{itemize}
      \item A forma em taxa erra $O(\Delta p_c^2)$ por passo.
      \item Com $S_a^{n+1}$ no lugar de $S_a^n$, o balanço do ar erra 0,25 a 0,34\,\%.
      \item Com a forma secante, a água fecha com $2\times10^{-14}$ e o ar com $8\times10^{-16}$.
    \end{itemize}
  \item \textbf{Newton completo}, com Jacobiano consistente:
    \begin{itemize}
      \item busca linear (o passo é dividido por 2 até 4 vezes);
      \item convergência quando $\|\bm R_u\|/\|\bm F_g\| < 10^{-7}$ e $\max|\delta p| < 10^{-6}$\,kPa, a partir da 2.ª
        iteração;
      \item no máximo 40 iterações (25 com as pressões congeladas);
      \item se falha, o passo é refeito com $\Delta t/2$.
    \end{itemize}
  \item \textbf{Não simetria:} o Jacobiano não é simétrico, pelo sinal das linhas de massa, por $S_\pi \neq \chi_\pi$ e
    pelos termos advectivos. A LU é não simétrica (PARDISO, ou skyline). No FS só entra $\bm J_{uu}$, que é simétrico.
  \item \textbf{Estrutura, sem plasticidade} (método $\theta$ com $\theta = 1$):
    \[
      \begin{bmatrix} \bm K & -\bm Q_{uw} & -\bm Q_{ua} \\ \bm Q_{wu} & \bm C_{ww} + \Delta t\,\bm H_w & \bm C_{wa} \\
      \bm Q_{au} & \bm C_{aw} & \bm C_{aa} + \Delta t\,\bm H_a \end{bmatrix}
      \begin{bmatrix} \delta\bm u \\ \delta\bm p_w \\ \delta\bm p_a \end{bmatrix} = -\bm R .
    \]
  \item \textbf{Contorno:} a percolação e a chuva usam uma condição unilateral. Cada vértice alterna entre ``de fluxo'' e
    ``alagado'' até o conjunto não mudar.
  \item \textbf{Estado inicial:} carregamento geostático com as pressões congeladas; depois $\bm u$ é zerado e
    $\bm\sigma'$ é mantido.
\end{itemize}

\paragraph{Casos particulares.}
\begin{itemize}
  \item saturado (Biot), com $1/M = 1/N + n/K_w$;
  \item ar passivo, $p_a = 0$;
  \item desacoplado, como em Rahardjo: Richards com o esqueleto rígido e acoplamento de uma via.
\end{itemize}

\paragraph{Decisões aplicadas no código.}
\begin{itemize}
  \item $\chi = S_w$ é o padrão;
  \item o armazenamento do ar usa $S_a^n$;
  \item só o PV2 com $\psi = \phi'$.
\end{itemize}
Depois dos documentos, o código ganhou:
\begin{itemize}
  \item o multiplicador de condutividade e a resistência por ponto, para os campos aleatórios;
  \item o PARDISO;
  \item o FS por arc-length;
  \item a chuva do \texttt{Common} (\texttt{SlopeRain}), com o esqueleto rígido e o ar passivo.
\end{itemize}

\paragraph{Pendências de documentação.}
\begin{itemize}
  \item As linhas do código citadas nos dois documentos estão defasadas.
  \item O \texttt{formulacao.tex} diz ``área'' onde o código usa o comprimento de influência.
  \item A classe Mohr--Coulomb antiga ainda é instanciada no material, embora nenhum projeto a use.
  \item O teste de Taylor não cobre $\chi = S_e$.
  \item Em aberto: D16 (estado inicial do talude), D17 (o ar na superfície alagada) e D19 (o critério de convergência
    sem a norma da massa).
\end{itemize}

\section{Verificação do modelo}
\label{sec:verificacao}

\paragraph{A ordem.} A revisão bibliográfica (31/10, 20 páginas, 35 trabalhos) conclui que nenhuma referência traz
$\bm u$, $p_w$ e $p_a$ medidos junto com uma solução analítica do problema não linear com Mohr--Coulomb. Por isso a
verificação vai por etapas: Richards, Biot saturado, Liakopoulos com ar passivo, dois fluidos, centrífuga, e só depois
as aplicações. O fator de segurança e a variabilidade espacial vêm depois de verificar a conservação de massa, a chuva
e a resposta constitutiva.

\paragraph{Resultados} (\path{Projects/README.md} e os README dos projetos; todos os testes rápidos passam).

\begin{center}\small
\begin{longtable}{@{}L{3.8cm} L{4.8cm} L{7.0cm}@{}}
\toprule
Projeto & O que verifica & Resultado \\
\midrule
\endhead
\texttt{UnsatMaterialPoint} & leis hidráulicas, Bishop, Mohr--Coulomb PV2, tangente &
  408 verificações; PV2 contra o artigo e contra a projeção KKT em 20\,000 estados; tangente de ordem 2 \\
\texttt{UnsatWaterPhase} & água com ar passivo e esqueleto rígido &
  hidrostático $1{,}8\times10^{-15}$\,kPa; Gardner 0,0015\,kPa; balanço $2\times10^{-14}$; infiltração de Philip a
  2,0\,mm da solução semianalítica \\
\texttt{UnsatAirPhase} & água e ar com esqueleto rígido &
  Boyle $8\times10^{-16}$; o limite de ar passivo converge com a mobilidade do ar; ar livre $5\times10^{-16}$ \\
\texttt{UnsatCoupled} & o modelo completo &
  Jacobiano de ordem 2,000 com 96 pontos plásticos; Terzaghi $\leq 0{,}0028\,p_0$; Mandel até $0{,}0038\,p_0$;
  Liakopoulos com ar passivo contra o experimento: RMS da vazão de $0{,}225\times10^{-6}$\,m/s e de $p_w$ de 0,44\,kPa,
  contra 0,227 e 0,48 da solução de Callari e Abati \\
\texttt{UnsatSRM}, \texttt{UnsatGIM} & talude seco contra Bishop &
  FS$_{SRM} = 1{,}2109$ e FS$_{GIM} = 1{,}8146$, contra 1,206 e 1,806 (malha nf $= 80$) \\
\texttt{LearningLiakopoulos} & a coluna contra uma referência 1D e o artigo &
  iguais em todos os dígitos impressos à referência 1D. A vazão tardia fica 10\,\% abaixo das curvas calculadas de
  Callari; com o passo de tempo do artigo, a referência reproduz o artigo, então a diferença vem do passo \\
\bottomrule
\end{longtable}
\end{center}

\paragraph{A coluna de Liakopoulos} (Callari e Abati, 2009).
\begin{itemize}
  \item \textbf{O ensaio:} areia de 1\,m de altura. Antes de $t = 0$ há fluxo permanente para baixo com $p_w = 0$ em
    toda a coluna, o que não é um estado hidrostático. Em $t = 0$ a alimentação é cortada, e a coluna drena pela base.
  \item \textbf{Os parâmetros:} $E = 1300$\,kPa, $\nu = 0{,}4$, $n = 0{,}2975$, $k_s = 4{,}41\times10^{-6}$\,m/s,
    $S_w = 1 - c_L p_c^{b_L}$ e $k_{rw} = 1 - 2{,}207(1 - S_w)^{1{,}0121}$, válidas para $S_w > 0{,}91$.
  \item \textbf{O que domina:} o armazenamento é dominado pela retenção,
    $nC_s \approx 2{,}7\times10^{-3} \gg 1/M = 3{,}6\times10^{-4} \gg n/K_w$ (kPa$^{-1}$). Por isso a coluna parcialmente
    saturada drena muito mais devagar que a saturada: a saturada recalca 1,761\,mm, com 90\,\% em cerca de 680\,s; a
    parcial ainda está em 1,585\,mm aos 7500\,s.
  \item \textbf{Diferenças em relação ao artigo:} o artigo usa a pressão de Coussy, o armazenamento em taxa e elementos
    P1/P0. O código usa Bishop com $\chi = S_w$, o armazenamento secante (conservativo) e Taylor--Hood P2--P1.
  \item \textbf{O que $p_a = 0$ não testa:} o transporte do ar. O ar ativo reduz a vazão (2,51 contra
    $3{,}38\times10^{-6}$\,m/s em 300\,s), mas não há $p_a$ medido para validá-lo.
\end{itemize}

\paragraph{Cuidados que a revisão aponta.}
\begin{itemize}
  \item \textbf{$\chi$ muda o problema.} Rahardjo equivale a $\chi = 1$, e o $\chi = S_w$ do material já é outro caso.
    O \texttt{SlopeRandomField} usa $\chi = 1$, como o Ceron e o Rahardjo.
  \item \textbf{Os métodos não são intercambiáveis.} SRM, GIM, equilíbrio limite e limite superior medem coisas
    diferentes. Uma falha da Newton não prova ruptura física.
  \item \textbf{Erros nas referências.} O Mandel de Haghighat et al.\ usa $K_{dr}$ no lugar de $K_{dr} + 4G/3$ (com ele,
    o erro chega a $0{,}23\,p_0$), e a Eq.\ 9 de Li et al.\ (Gardner) tem o sinal de $Q$ trocado.
  \item \textbf{Critérios de concordância:} três malhas e três passos de tempo, erros $L_2$ e máximo, a tolerância da
    Newton bem abaixo do erro de discretização, e o balanço de massa por fase. A massa de ar é
    $\int \rho_a n S_a\, dV$.
\end{itemize}

\paragraph{Pendências da verificação.}
\begin{itemize}
  \item A norma do resíduo de massa ainda não entra no critério de convergência.
  \item $p_a$ fica mal definida quando $S_a \to 0$.
  \item Ficam em aberto o estado inicial do talude (o teto de sucção de 75\,kPa) e a condição do ar na superfície
    alagada.
  \item Bergamaschi--Putti, Barry--Mercer, Xie et al.\ e a centrífuga de Wang et al.\ ainda não foram implementados.
  \item Há divergências de documentação:
    \begin{itemize}
      \item a revisão descreve a versão anterior ($\chi = 1$ como padrão);
      \item os resumos apresentam van Genuchten--Mualem, mas o padrão do \texttt{TPZUnsaturatedHydraulics} é
        Fredlund--Xing com $k_{rw} = S_e^4$; os testes fixam van Genuchten explicitamente.
    \end{itemize}
\end{itemize}

\section{Fator de segurança: SRM, GIM e arc-length}
\label{sec:fs}

Fonte: o relatório \emph{Cálculo do fator de segurança de um talude não saturado sob chuva: SRM e GIM, com e sem
arc-length}, de 7/11, com 34 páginas (\path{arclength/relatorio_fs}, igual a \path{docs/fs_arc_length}). O
objetivo é escolher o método e os parâmetros do FS que o Monte Carlo vai calcular milhares de vezes. Os três
critérios: o método tem de responder sempre, ficar perto do valor verdadeiro e custar pouco.

\paragraph{Definições.} Em ambos, as pressões ficam congeladas: o estado é salvo, $p_w$ e $p_a$ são fixados, e só o
equilíbrio é resolvido, a partir de tensão nula.
\begin{itemize}
  \item \textbf{SRM:} o fator $F$ divide $c'$, $\tan\phi'$ e $\tan\psi$, e só $\bm f_{int}$ depende dele.
  \item \textbf{GIM:} o fator $\lambda$ multiplica o peso e a parte positiva de $p_B$. A sucção fica inteira, porque é
    resistência, e a resistência é a real. Só $\bm f_{ext}$ depende de $\lambda$, e linearmente.
  \item \textbf{O que é o FS:} o maior fator para o qual existe equilíbrio.
    \begin{itemize}
      \item O SRM e o GIM medem coisas diferentes: os valores não se comparam, só o lado de 1 em que ficam.
      \item O $\Gamma$ do Ceron é o fator do GIM. O README do \texttt{SlopeRandomField} o escreve como
        $\Gamma = P_{mr}/(P_\gamma + P_u)$.
    \end{itemize}
  \item \textbf{Hipótese de trabalho} (Sysala et al., 2025): com fluxo associado e sem história plástica, o fator cresce
    monotonamente ao longo do caminho. Então todo equilíbrio convergido é uma cota inferior do FS.
\end{itemize}

\paragraph{O problema do método antigo.} O método antigo usava uma lista de fatores com bisseção e tomava a falha da
Newton como ruptura. Uma falha não distingue ``não há equilíbrio'' de ``o salto foi grande demais''. No GIM em 24\,h,
isso deu 0,992 em vez de 1,154, um valor abaixo de 1 que contraria a cota inferior.

\paragraph{Arc-length.} O fator $\Lambda$ passa a ser incógnita, e o comprimento do passo no espaço $(\bm u, \Lambda)$ é
prescrito. O método traça o trecho íngreme, o joelho e o patamar, e o FS é o maior $\Lambda$ do caminho.
\begin{gather}
  \bm r(\bm u, \Lambda) = \bm f_{int}(\bm u, \Lambda) - \bm f_{ext}(\Lambda), \qquad
  \bm K_T\,\delta\bm u - \bar{\bm f}\,\delta\Lambda = -\bm r, \qquad
  \delta\bm u = \delta\bm u^* + \delta\Lambda\,\delta\bar{\bm u}, \\
  \bm K_T\,\delta\bm u^* = -\bm r, \qquad \bm K_T\,\delta\bar{\bm u} = \bar{\bm f}, \qquad
  \Delta\bm u\cdot\Delta\bm u + \alpha_e^2\,\Delta\Lambda^2 = l^2 .
\end{gather}
\begin{itemize}
  \item \textbf{Vetor de referência:} no GIM, $\bar{\bm f} = \bm r(\bm u, 0) - \bm r(\bm u, 1)$, constante. No SRM,
    $\bar{\bm f}$ é calculado por diferença finita em $F$, e só os pontos plásticos contribuem.
  \item \textbf{Termo $\alpha$:} sem ele, o SRM no trecho elástico tem $\bar{\bm f} = \bm 0$ e o passo não controla $F$.
    O termo $\alpha_e^2 \Delta\Lambda^2$ dá $\Delta F = l/\alpha_e$ e dispensa a lista.
  \item \textbf{Corretor:} uma equação do segundo grau em $\delta\Lambda$. Fica a raiz de menor ângulo com o incremento;
    sem raiz real, $l$ é cortado.
  \item \textbf{Preditores:} tangente, com a superfície de escoamento encolhida em $10^{-6}$ (é essencial), ou secante.
  \item \textbf{Comprimento adaptativo:} $l_{n+1} = l_n \min(2, \max(0{,}5, \sqrt{10/\mathrm{iter}_n}))$.
  \item \textbf{Parada:} $\Delta\Lambda\, l_{pedido}/l < \mathrm{tol}$. Assim, um corte de $l$ não se confunde com o
    patamar.
\end{itemize}

\paragraph{Resultados.} São 43 casos, mais de 1600 rodadas, e 38 casos nas tabelas finais. O erro é medido contra o
maior equilíbrio achado por qualquer método.
\begin{itemize}
  \item \textbf{Configuração final:}
    \begin{itemize}
      \item o pior erro é de $-0{,}46\,\%$ no SRM e $-0{,}41\,\%$ no GIM;
      \item contra o método antigo, as iterações caem 35\,\% no SRM e 7\,\% no GIM;
      \item o GIM antigo chegava a $-26{,}9\,\%$, com 18 de 35 casos acima de 1\,\%;
      \item com a tolerância de 0,01, o SRM chega a $-4{,}85\,\%$.
    \end{itemize}
  \item \textbf{FS ao longo do tempo no caso base} ($H = 10$\,m, $45^\circ$, solo f50, chuva de 36\,mm/h por 24\,h):
    \begin{itemize}
      \item no SRM, de 2,116 em 0\,h a 1,060 em 24\,h, voltando a 1,280 em 48\,h;
      \item no GIM, de 5,696 a 1,153 e depois a 2,010.
    \end{itemize}
  \item \textbf{Modelo acoplado com ar ativo:}
    \begin{itemize}
      \item o mínimo em 24\,h é de 1,005 no SRM e 1,018 no GIM, ambos $\geq 1$, como a cota inferior prevê;
      \item o ar faz o FS cair mais cedo: 1,110 contra 1,426 em 12\,h;
      \item esse efeito depende do piso de $k_{ra}$, então FS $< 1$ nesse modelo não é conclusão firme.
    \end{itemize}
\end{itemize}

\paragraph{Parâmetros recomendados} (\texttt{ParamsSRM} e \texttt{ParamsGIM}).
\begin{itemize}
  \item \textbf{Comuns:} tol $= 0{,}001$, preditor secante, superfície encolhida em $10^{-6}$, alvo de 10 iterações,
    maxIt $= 12$, fatorMax $= 10$.
  \item \textbf{SRM:} $\alpha = 0{,}1$, sem lista, $l_0 = 0{,}01$, $F_0 = 0{,}6$, com reinício em 0,4.
  \item \textbf{GIM:} $\alpha = 0$, $\lambda_0 = 0{,}5$, $l_0 = 0{,}25$, $l_{max} = 2$.
\end{itemize}

\paragraph{No código.}
\begin{itemize}
  \item \textbf{Arc-length:} \texttt{TPZUnsatAnalysis} (\texttt{ArcLengthStep}, \texttt{ArcLengthStepSRM},
    \texttt{ArcLengthIterate}).
  \item \textbf{\texttt{SlopeRandomField}:} \path{FATOR=gim_arco} usa o GIM com arc-length, mas o padrão continua
    sendo a lista do SRM.
  \item \textbf{\path{Projects/Common/SlopeSolver}:} \texttt{ComputeFactor} oferece os dois métodos, com e sem
    arc-length, e é usado pelo \texttt{SlopeMonteCarlo}.
  \item \textbf{\texttt{UnsatSRM} e \texttt{UnsatGIM}:} testam o talude seco contra Bishop (FS$_{SRM} = 1{,}2109$ e
    FS$_{GIM} = 1{,}8146$, contra 1,206 e 1,806), com lista e bisseção.
\end{itemize}

\paragraph{Divergências entre o relatório e o código.}
\begin{itemize}
  \item O \texttt{calculaFS} do \texttt{LearningMohrCoulomb2} usa tol $= 0{,}01$, e não 0,001.
  \item O \texttt{SRMArcLength} do \texttt{Common} não segue o \texttt{ParamsSRM}: parte de $F = 1$ e não consegue dar
    FS $< 1$.
  \item O código do estudo de sensibilidade não está no repositório.
\end{itemize}

\section{A chuva: o talude de Rahardjo et al.\ (2007)}
\label{sec:chuva}

A linha da chuva tem três produtos.
\begin{enumerate}
  \item \textbf{A reprodução determinística do Rahardjo} (notas \texttt{ar4} e \path{secoes_rahardjo}, 28/10; o
    relatório \emph{Contraexemplos a Rahardjo et al.\ (2007)}, 29/10, em \path{docs/biot_nao_saturado}).
  \item \textbf{O artigo em preparação}, \emph{Critical rainfall intensity of unsaturated slopes}
    (\path{docs/biot_nao_saturado/artigo}). Ele mede quanto as hipóteses usuais (esqueleto rígido, ar passivo,
    desacoplado) superestimam a chuva que leva à ruptura.
  \item \textbf{O relatório da chuva} (\path{relatorio_form/relatorio_chuva}, 14/11, 12 páginas): o Monte Carlo da
    chuva com o método estocástico do Ceron. Está só no zip.
\end{enumerate}

\paragraph{O modelo do relatório.}
\begin{itemize}
  \item \textbf{Geometria:} domínio de 70\,m; o talude tem $H = 10$\,m e $45^\circ$. O nível d'água fica 5\,m abaixo
    do pé e sobe $7^\circ$ para dentro do maciço. Acima dele a pressão é hidrostática, com a sucção limitada a 75\,kPa.
  \item \textbf{Solo:} $c' = 10$\,kPa, $\phi' = \phi^b = 26^\circ$, $\gamma = 20$\,kN/m$^3$, com três versões: f10
    ($a = 10$\,kPa e $k_s = 360$\,mm/h), f50 (50 e 36\,mm/h) e f100 (100 e 3,6\,mm/h).
  \item \textbf{Hidráulica:} Fredlund--Xing com $C(\psi) = 1$ e $k_w = k_s\,\Theta^4$:
    \[
      \Theta = \theta_w/\theta_s = \big[\ln\big(e + (\psi/a)^n\big)\big]^{-m}, \qquad \psi = u_a - u_w = -p_w .
    \]
  \item \textbf{Transiente}, com as hipóteses do SEEP/W: esqueleto rígido e ar passivo, isto é, Richards em $p_w$, o
    material com o acoplamento desligado. Euler implícito com $\Delta t = 0{,}1$\,h.
  \item \textbf{Chuva:} intensidade $I$ por 24\,h. Em cada nó da superfície vale uma de duas condições: $q_{in} = I$ com
    $p_w \leq 0$, ou $p_w = 0$ com $q_{in} \leq I$ (face de percolação). O conjunto de nós da face é iterado dentro do
    passo até se repetir.
  \item \textbf{Fator de segurança:}
    \begin{itemize}
      \item em cada instante, $p_w$ é congelada e interpolada para uma malha do fator, refinada em 4 níveis na zona
        plástica;
      \item o SRM usa $\chi = 1$, com $\tau_f = c' + (\sigma_n - u_a)\tan\phi' + (u_a - u_w)\tan\phi^b$;
      \item o caminho é o peso de 0 a 1, seguido de arc-length em $F$.
    \end{itemize}
  \item \textbf{Campos:} $c'$, $\phi'$ e $k_s$ lognormais (CoV 0,3, 0,1 e 0,6), KL analítica no retângulo de
    $70 \times 30$\,m com $M = 2000$, que retém 96,4\,\% da variância.
\end{itemize}

\paragraph{Resultados do relatório.}
\begin{itemize}
  \item \textbf{Validação contra a Fig.\ 5 do Rahardjo:}
    \begin{itemize}
      \item FS$(0) = 2{,}0975$, 1\,\% abaixo dos 2,118 do artigo;
      \item a diferença média é de $+1{,}6\,\%$ no f10, $0{,}0\,\%$ no f50 e $-0{,}2\,\%$ no f100;
      \item o f50 chega ao mínimo de 1,0346 em 24\,h (o artigo dá 1,03);
      \item o f10 cruza 1 em 5,85\,h.
    \end{itemize}
    Durante a chuva o FS só cai, e por isso a campanha usa FS(24\,h).
  \item \textbf{GIM contra SRM:} antes da chuva $\Gamma \approx 2{,}5\,$FS (5,28 contra 2,10), porque a sucção fica
    inteira no GIM. Em 24\,h os dois quase coincidem (1,088 contra 1,035). O relatório escolhe o SRM, que é menos
    ruidoso, embora custe cerca do dobro.
  \item \textbf{Intensidade de referência:} $I_{ref} = 31{,}05$\,mm/h dá $\Gamma = 1{,}3345$, o mesmo dos 1,336 do Ceron.
  \item \textbf{FS(24\,h) em função de $k$:} tem a forma de U, com o mínimo em $k \approx I$.
  \item \textbf{Pilotos} (8 a 10 realizações): o CoV de FS(24\,h) é de 5,6 a 7,6\,\% no SRM, contra 23,5\,\% do $\Gamma$ do
    Ceron. Nenhuma realização rompeu, e o ajuste lognormal sugere $P_f \approx 1\,\%$.
  \item \textbf{FORM com variáveis}, num solo homogêneo, sem campos: $\beta = 0{,}681$ e $P_f = 24{,}8\,\%$. Sem a média
    espacial, o fator espalha muito mais. É o limite de comprimento de correlação infinito.
  \item \textbf{Campanha:} 27 casos e 57\,000 realizações, cerca de 3,8 dias no fermi. Está pronta, mas não rodou.
\end{itemize}

\paragraph{O artigo: conclusões.} Elas valem para o talude do Rahardjo e os casos calculados.
\begin{itemize}
  \item \textbf{O critério de comparação.} Nos modelos acoplados com fluxo associado, o FS mínimo não cai abaixo de 1.
    Por isso os modelos são comparados pela intensidade crítica $I_{c,\delta}$, a que leva o FS mínimo a $1 + \delta$,
    e não pelo FS mínimo.
  \item \textbf{Quanto as hipóteses usuais erram.} Com $k_s = 36$\,mm/h:
    \begin{itemize}
      \item o modelo usual dá $I_{c,\delta} = 47$ a 49\,mm/h;
      \item o modelo completo dá 30 a 35\,mm/h com $\psi = \phi'$, e 26\,mm/h com $\psi = 0$;
      \item as hipóteses usuais superestimam a chuva crítica em 40 a 59\,\% e em 85 a 88\,\%, e atrasam a ruptura em
        4,5 e 9,9\,h.
    \end{itemize}
  \item \textbf{Ar e acoplamento.} O ar ativo reduz $I_{c,\delta}$ por um fator de 2 a 2,5; o acoplamento a aumenta. Sem
    dilatância, o transiente acoplado colapsa.
  \item \textbf{Robustez.} O viés persiste com $\chi = S_e$ e com a malha fina, e muda de sinal no solo permeável f10.
\end{itemize}

\paragraph{Pendências e cuidados.}
\begin{itemize}
  \item \textbf{O código do relatório da chuva não está no repositório.} Faltam \texttt{rahardjo.cpp},
    \texttt{montecarlo.cpp}, \texttt{RainTransient.h} e \path{form_variaveis.py}. A versão do repositório é anterior
    e dá FS 2,116 e 1,060, contra 2,0975 e 1,0346.
  \item \textbf{Escala dos campos.} $L_x = 20$\,m e $L_y = 2$\,m são do talude do Ceron, com $H = 5$\,m. Em relação a
    $H = 10$\,m, a referência equivale a metade da escala do Ceron, e o relatório compara $P_f$ com o Ceron sem
    comentar isso.
  \item \textbf{Chuva inclinada ou horizontal.} O relatório aplica a chuva por metro de superfície inclinada; as notas
    de 28/10 a aplicavam na projeção horizontal. Na face a $45^\circ$, a diferença é um fator $\sqrt 2$.
  \item \textbf{$I_{ref}$.} Foi calibrado com o $\Gamma$ do GIM, que não é o fator da campanha, e dentro do ruído do GIM.
  \item \textbf{Os mínimos depois da chuva.} O Monte Carlo para em 24\,h e não vê mínimos posteriores, como o do f100.
  \item \textbf{Os próximos passos do relatório:}
    \begin{itemize}
      \item rodar a campanha;
      \item com $P_f \approx 1\,\%$, aumentar $N$ (39\,600 realizações para um CoV de 5\,\%) ou usar a simulação por
        subconjuntos;
      \item o FORM com os campos, com a KL truncada e a amostragem por importância (Seção~\ref{sec:confiab});
      \item o ar ativo.
    \end{itemize}
\end{itemize}

\section{Campos aleatórios: Karhunen-Loève}
\label{sec:campos}

\paragraph{O campo.} Cada propriedade (por exemplo, $c'$) é lognormal:
\begin{equation}
  c'(\bm x) = \exp\big(\mu_{\ln} + \sigma_{\ln} Z(\bm x)\big), \qquad
  \sigma_{\ln}^2 = \ln(1 + \mathrm{CoV}^2), \quad \mu_{\ln} = \ln\mu - \tfrac12 \sigma_{\ln}^2 ,
\end{equation}
em que $Z$ é um campo gaussiano de média 0 e variância 1 com o núcleo separável do Ceron:
\begin{equation}
  C(\bm x, \bm x') = \exp\!\left(-\frac{|x - x'|}{L_x} - \frac{|y - y'|}{L_y}\right).
\end{equation}

\paragraph{A expansão de Karhunen-Loève (KL).}
\begin{equation}
  Z(\bm x) \approx \sum_{k=1}^{M} \sqrt{\lambda_k}\, \xi_k\, f_k(\bm x), \qquad
  \int_D C(\bm x, \bm x') f_k(\bm x')\, d\bm x' = \lambda_k f_k(\bm x), \qquad \xi_k \sim \mathcal N(0,1)
  \ \text{independentes}.
\end{equation}
\begin{itemize}
  \item Os autovalores somam a área: $\sum_k \lambda_k = |D|$.
  \item O erro médio de truncamento é $\varepsilon_M = 1 - \sum_{k\leq M} \lambda_k/|D|$, e a variância truncada em
    cada ponto é $\sigma_M^2(\bm x) = \sum_{k\leq M} \lambda_k f_k(\bm x)^2 < 1$.
  \item Os $\xi_k$ são as variáveis aleatórias do problema: $n = 3M$, ou seja, 6000 com $M = 2000$.
\end{itemize}

\paragraph{Três formas de obter os autopares.}
\begin{enumerate}
  \item \textbf{Forma fechada no retângulo que contém o domínio} (\texttt{KLSeparavel.h}; é a que o \texttt{SlopeMonteCarlo}
    usa hoje).
    \begin{itemize}
      \item O problema 1D com o núcleo $e^{-c|s - t|}$ em $[-a, a]$ tem modos pares, com $\omega \tan(\omega a) = c$, e
        ímpares, com $\omega + c\tan(\omega a) = 0$; o autovalor é $\lambda = 2c/(\omega^2 + c^2)$.
      \item Os modos 2D são os produtos $f_i(x)\,g_j(y)$ com $\lambda_i^x \lambda_j^y$, e ficam os $M$ maiores.
      \item O campo é avaliado ponto a ponto, nos pontos de integração.
    \end{itemize}
  \item \textbf{Galerkin} (NeoPZ, \texttt{TPZKarhunenLoeveMat}): $C\bm v = \lambda B \bm v$, com
    $B_{ij} = \int N_i N_j$ e $C_{ij} = \iint N_i(\bm x)\, C(\bm x, \bm x')\, N_j(\bm x')$. É a KL do próprio domínio, e a
    sua qualidade depende da malha e da ordem da integração dupla.
  \item \textbf{Nyström}: uma quadratura $(\bm x_j, w_j)$ do domínio, com a matriz simétrica $W^{1/2} K W^{1/2}$,
    $K_{ij} = C(\bm x_i, \bm x_j)$. A interpolação é
    $f_k(\bm x) = \lambda_k^{-1} \sum_j w_j\, C(\bm x, \bm x_j)\, f_k(\bm x_j)$.
    \begin{itemize}
      \item O seu caderno \texttt{nymstron.nb} faz a versão nodal, com os pesos de massa concentrada: a área de cada
        elemento dividida igualmente entre os seus nós, com $\sum w = |D|$.
      \item Sem os pesos, os autovalores de $K$ saem com a escala errada: $\mathrm{tr}\,K = n$ em vez de $|D|$.
    \end{itemize}
\end{enumerate}

\paragraph{O que se concluiu.}
\begin{itemize}
  \item \textbf{Autovalores do talude do Ceron:} $\lambda_1 = 48{,}81$, convergido pelo Galerkin, pelo Nyström e por uma
    referência extrapolada.
  \item \textbf{O retângulo restrito ao talude tem a covariância certa.} Quando $M \to \infty$, a forma fechada do
    retângulo, restrita ao talude, tem a covariância pedida: o campo do talude é a restrição de um campo estacionário.
    \begin{itemize}
      \item Os autovalores são outros, porque a base é a do retângulo, mas a estatística no talude é a mesma.
      \item Com $M = 2000$, $\varepsilon_M = 0{,}0079$, em 1,3\,s; o Galerkin no talude parou em cerca de 0,029, limitado
        pela malha, em 260\,s.
      \item Por termo, a base do retângulo é menos eficiente: 63 termos para igualar 50 do talude.
    \end{itemize}
  \item \textbf{Mesmos $\xi_k$ e a mesma realização} (\texttt{VerifyRandomFields}, branch próprio: coesão nos 21\,069
    pontos de integração de uma malha uniforme de 0,25\,m):
    \begin{itemize}
      \item \emph{Galerkin $\times$ Nyström}: é a mesma realização. A diferença RMS é de 0,0018\,kPa com 10 termos e
        0,0055\,kPa com 20, para um desvio do campo de cerca de 3\,kPa. Com 200 termos, sobe para 0,59\,kPa, porque
        nenhuma das duas resolve os modos de ordem acima de 25 a 40.
      \item \emph{Galerkin $\times$ forma fechada}: a diferença RMS é de 3,85\,kPa, com correlação 0,49. É outra
        realização da mesma estatística, não um erro.
      \item \emph{Estatística igual nos três}: com 200 termos, o déficit de variância é de 0,042, 0,039 e 0,044 (o exato
        é 0,0375). A covariância difere do núcleo em $10^{-3}$ a $2\times 10^{-3}$, e a média e o desvio de Monte Carlo
        ficam dentro do ruído amostral.
    \end{itemize}
  \item \textbf{GSTools} (método da randomização): com o espectro produto de duas Cauchy, ele reproduz o núcleo
    separável. O \texttt{gs.Exponential} anisotrópico usa a distância euclidiana escalada, que é outro núcleo.
  \item \textbf{Mika et al.\ (2021)} (\texttt{KLCilindroMika}): o Galerkin do NeoPZ com a geometria exata do meio
    cilindro converge com taxa $h^4$. Os autovalores do caderno \texttt{randomfied-v2.nb} não se reproduzem.
  \item \textbf{A escolha que está em aberto:}
    \begin{itemize}
      \item você prefere o Nyström, que é a KL do próprio domínio;
      \item a forma fechada é estatisticamente equivalente e mais barata;
      \item a escolha muda a realização, não a estatística;
      \item o Nyström precisa de menos termos para a mesma variância, e isso importa para o FORM e para metamodelos.
    \end{itemize}
\end{itemize}

\section{Monte Carlo e refinamento por realização}
\label{sec:mc}

\paragraph{O estimador.}
\begin{equation}
  \hat P_f = \frac{n_f}{S}, \qquad \mathrm{CoV}(\hat P_f) = \sqrt{\frac{1 - P_f}{S\, P_f}} .
\end{equation}
Com $P_f \approx 11\,\%$, 1000 realizações dão um CoV de 9\,\%, e 10\,000 dão 2,8\,\%. O ajuste lognormal de $\Gamma$ dá
uma segunda estimativa, $P(\Gamma < 1)$, que depende do modelo de distribuição.

\paragraph{Refinamento por realização} (\texttt{SlopeRandomFieldRefina}, relatório de 11/11).
\begin{enumerate}
  \item \textbf{Etapa 1:} GIM na malha grossa.
  \item \textbf{Etapa 2:} a banda de ruptura, $\sqrt{J_2(\bm\varepsilon_p)} \geq 0{,}1 \max$, define um campo de tamanhos
    para o \texttt{gmsh}.
  \item \textbf{Etapa 3:} GIM na malha refinada.
\end{enumerate}
O contorno é o mesmo, então o retângulo da KL e o campo também são os mesmos. Os resultados:
\begin{itemize}
  \item \textbf{Malha recomendada:} a \texttt{foco375} mais um passo até $h = 0{,}125$\,m. A diferença média em
    $\Gamma$ contra a \texttt{foco125} é de $+0{,}14\,\%$, e o viés em $P_f$, de $-1\,\%$, fica abaixo do erro de Monte
    Carlo. Os passos até 0,25 e 0,1875\,m subestimam $P_f$ em 17 e 10\,\%.
  \item \textbf{Malha inicial mais grossa} (\texttt{g15}): dá os mesmos erros por nível, mas a banda é mais larga, e só
    compensa nos níveis grossos.
  \item \textbf{A \texttt{foco125} não está convergida:} a extrapolação indica $\Gamma$ cerca de 1\,\% alto, o que
    deslocaria $P_f$ em cerca de $+9\,\%$.
\end{itemize}

\paragraph{O pipeline atual} (\path{Projects/Common} e \texttt{SlopeMonteCarlo}, commits de 11 e 12/11).
\begin{itemize}
  \item \textbf{Em cada realização:} percolação permanente com o $k_v$ da realização, depois $\Gamma$ pelo GIM com
    arc-length e as pressões congeladas, com refinamento adaptativo hp da zona plástica em 4 níveis.
  \item \textbf{Campanha:} 10\,000 realizações por caso, para os 26 casos das Seções 6.1 a 6.5 do artigo, com
    continuação automática dos jobs no fermi; cerca de 95\,s por tarefa.
  \item \textbf{Saídas:} o $\Gamma$ de cada nível, a banda e a massa que desliza, com $c'$, $\phi'$ e $k$ médios para o
    Spearman da Seção 6.1.
\end{itemize}

\section{Confiabilidade: FORM, Monte Carlo e amostragem por importância}
\label{sec:confiab}

Fonte: \path{Projects/Confiabilidade} (relatório \texttt{confiabilidade.pdf} e código \texttt{confiabilidade.wl}).

\paragraph{Espaços e transformação.}
\begin{itemize}
  \item $\bm X$, as variáveis físicas, vão para $\bm Z$, normais correlacionadas, por $z_i = \Phi^{-1}(F_i(x_i))$.
  \item $\bm Z$ vai para $\bm U$, normais padrão independentes, por $\bm u = L^{-1}\bm z$, com $R_Z = LL^T$.
  \item Com marginais lognormais, $\rho_Z = \ln(1 + \rho\,\delta_i\delta_j)/(\zeta_i\zeta_j)$ (Nataf).
  \item No talude com KL, os $\xi_k$ já são o espaço $\bm U$, e $G(\bm u) = \Gamma - 1$.
\end{itemize}

\paragraph{FORM.} O ponto de projeto é $\bm u^* = \arg\min\{\|\bm u\| : G(\bm u) = 0\}$, com $\beta = \|\bm u^*\|$ e
$P_f \approx \Phi(-\beta)$. Os fatores $\alpha_i^2$ dizem quanto cada variável pesa.
\begin{itemize}
  \item O algoritmo é o HL-RF:
    $\bm u_{k+1} = \dfrac{\nabla G_k\cdot\bm u_k - G_k}{\|\nabla G_k\|^2}\,\nabla G_k$. É preciso a busca de Armijo, ou o
    método oscila nos casos secos.
  \item O FORM aproxima a superfície pelo plano tangente: o erro depende da curvatura e não é conhecido.
\end{itemize}

\paragraph{Amostragem por importância.} A densidade é $h = \varphi_n(\bm v - \bm u^*)$, o peso
$w = \exp(-\bm v\cdot\bm u^* + \tfrac12\|\bm u^*\|^2)$, e $\hat P_f = \frac1N\sum I[G \leq 0]\,w$, com o CoV da amostra.
\begin{itemize}
  \item A estimativa não tem viés.
  \item A eficiência não piora com a dimensão, desde que o deslocamento esteja certo e a covariância não mude.
\end{itemize}

\paragraph{Resultados.}
\begin{itemize}
  \item \textbf{Seu código original} ($g = X_1X_2 - X_3$):
    \begin{itemize}
      \item o $\beta = 3{,}049$ estava certo;
      \item mas $L$ e $L^{-1}$ estavam trocados na transformação, e a volta $z \to x$ estava errada. Isso só aparece
        com correlação;
      \item com $P_f = 1{,}18\times10^{-3}$, Monte Carlo com 1000 amostras tem CoV de 92\,\%. A amostragem por
        importância com 1000 amostras tem CoV de 6\,\%.
    \end{itemize}
  \item \textbf{Talude infinito sob chuva} (parâmetros do Ceron, carga de pressão $h$ como a variável da chuva):
    \begin{itemize}
      \item $P_f$ vai de $10^{-12}$ a 0,6 ao longo dos cenários;
      \item o FORM superestima $P_f$ em 6 a 42\,\%;
      \item a carga de pressão domina ($\alpha_h^2 \approx 0{,}65$ a $0{,}69$);
      \item a correlação negativa entre $c'$ e $\phi'$ reduz $P_f$.
    \end{itemize}
  \item \textbf{No talude de elementos finitos:}
    \begin{itemize}
      \item com $P_f \approx 11\,\%$, Monte Carlo basta;
      \item abaixo de cerca de $10^{-2}$, use amostragem por importância, com o deslocamento igual à média dos $\bm\xi$ das
        realizações que falham, restrita aos primeiros modos;
      \item o FORM direto não é viável: o gradiente custa $2n$ análises, com $n = 6000$.
    \end{itemize}
\end{itemize}

\section{O que está pronto e o que falta}
\label{sec:pendencias}

\paragraph{Pronto e verificado.}
\begin{itemize}
  \item O modelo $\bm u$--$p_w$--$p_a$ com o PV2:
    \begin{itemize}
      \item Jacobiano de ordem 2;
      \item conservação de massa exata;
      \item Terzaghi, Mandel, Philip e Gardner;
      \item Liakopoulos com ar passivo;
      \item talude seco a 0,5\,\% do Bishop.
    \end{itemize}
  \item O FS por arc-length, com os parâmetros escolhidos em 43 casos (erro $\leq 0{,}5\,\%$).
  \item A reprodução determinística do Rahardjo, a cerca de 2\,\%.
  \item As três formas da KL, comparadas entre si (\texttt{VerifyRandomFields}).
  \item O refinamento por realização e a escolha da malha.
  \item O Monte Carlo do Ceron com 3 níveis, que dá $P_f = 10{,}3\,\%$ contra 11,5\,\%.
\end{itemize}

\paragraph{Para rodar.}
\begin{itemize}
  \item A campanha do Ceron: 26 casos $\times$ 10\,000 realizações com 4 níveis, cerca de 6 dias com 4 nós.
  \item A campanha da chuva: 27 casos, 57\,000 realizações, cerca de 3,8 dias, com o código que precisa ir para o
    repositório antes.
\end{itemize}

\paragraph{Pendências, por prioridade.}
\begin{enumerate}
  \item \textbf{Colocar no repositório o código do relatório da chuva}: \texttt{rahardjo.cpp}, \texttt{montecarlo.cpp},
    \texttt{RainTransient.h}, \path{form_variaveis.py} e \path{cluster/}. Sem ele, os resultados não são
    reproduzíveis, e a versão do repositório é outra.
  \item \textbf{Alinhar as tolerâncias do FS}:
    \begin{itemize}
      \item o \texttt{calculaFS} do \texttt{LearningMohrCoulomb2} usa tol $= 0{,}01$, e o recomendado é 0,001;
      \item o \texttt{SRMArcLength} do \texttt{Common} não segue o \texttt{ParamsSRM} e não dá FS $< 1$;
      \item o padrão do \texttt{SlopeRandomField} ainda é a lista do SRM.
    \end{itemize}
  \item \textbf{Chuva com $P_f \approx 1\,\%$}:
    \begin{itemize}
      \item o Monte Carlo precisa de cerca de $10^4$ a $4\times10^4$ realizações por caso;
      \item ou amostragem por importância com o deslocamento das realizações que falham, ou a simulação por
        subconjuntos. O \texttt{RandomFields::Sample} recebe a mudança.
    \end{itemize}
  \item \textbf{A escala dos campos na chuva} ($L_x$, $L_y$ contra $H = 10$\,m) e a chuva na superfície inclinada ou na
    projeção horizontal: decidir e documentar.
  \item \textbf{A diferença para o artigo do Ceron}: $\Gamma$ de 1,373 contra 1,336 (o contorno distante da percolação)
    e $P_f$ de 10,3 contra 11,5\,\%.
  \item \textbf{A KL}: decidir entre a forma fechada (a atual) e o Nyström (a sua preferência). A estatística é a mesma;
    o Nyström precisa de menos termos.
  \item \textbf{Em aberto na formulação}:
    \begin{itemize}
      \item D16: o estado inicial do talude, com o teto de sucção;
      \item D17: o ar na superfície alagada;
      \item D19: a norma da massa no critério de convergência;
      \item $p_a$ quando $S_a \to 0$;
      \item o teste de Taylor com $\chi = S_e$;
      \item a classe Mohr--Coulomb antiga, ainda instanciada.
    \end{itemize}
  \item \textbf{Validação do ar ativo}: só qualitativa. A revisão sugere Touma e Vauclin (1986).
  \item \textbf{Documentação}:
    \begin{itemize}
      \item as linhas do código citadas estão defasadas;
      \item a revisão descreve a versão anterior ($\chi = 1$ como padrão);
      \item os resumos apresentam van Genuchten--Mualem, e o padrão do código é Fredlund--Xing;
      \item a notação muda entre os documentos.
    \end{itemize}
\end{enumerate}

\section{Roteiro de estudo}
\label{sec:roteiro}

\paragraph{A ordem de leitura.} Para cada etapa, o que dominar.
\begin{enumerate}
  \item \textbf{Revisão bibliográfica} (20 p.). Por que verificar por etapas; por que $\chi$ muda o problema; por que
    SRM e GIM não são intercambiáveis.
  \item \textbf{Resumo do não saturado} (13 p.), depois \textbf{\texttt{formulacao}} (25 p.), ou o \texttt{verick} se
    preferir a ordem da tese de Santos. A dominar: as equações de campo, o armazenamento secante, os resíduos e a
    estrutura do Jacobiano.
  \item \textbf{Liakopoulos}: o resumo de 24 p., Seções 10 e 11, e a tabela de verificação do
    \path{Projects/README.md}. A dominar: o que cada teste verifica e o que não verifica.
  \item \textbf{\path{relatorio_fs}} (34 p.). As definições do SRM e do GIM, o arc-length, os critérios de parada e
    os parâmetros recomendados.
  \item \textbf{A chuva}: o \path{relatorio_chuva} (12 p.) e as conclusões do artigo. O transiente, a condição de
    superfície, FS($t$) e a intensidade crítica.
  \item \textbf{Campos aleatórios}: o relatório de 11/11 (Seções 4 a 6), o README do \texttt{VerifyRandomFields} e o
    seu \texttt{nymstron.nb}.
  \item \textbf{Monte Carlo}: o relatório de 11/11 (Seções 2 e 3) e o README do \path{SlopeMonteCarlo/cluster}.
  \item \textbf{Confiabilidade}: o \texttt{confiabilidade.pdf}.
\end{enumerate}

\paragraph{Perguntas para testar o entendimento}, com a resposta curta.
\begin{enumerate}[label=\textbf{P\arabic*.}, leftmargin=2.2em]
  \item \emph{Por que o armazenamento secante, $n\Delta S_w$, e não $n\Gamma\Delta p_c$?} Somando as equações nodais, ele
    dá a variação exata do volume armazenado (erro de $2\times10^{-14}$). A forma em taxa erra $O(\Delta p_c^2)$ por
    passo.
  \item \emph{Por que $S_a^n$ no termo do ar?} Porque $(m_a^{n+1} - m_a^n)/\rho_a^{n+1}$ fica exato. Com $S_a^{n+1}$
    sobra um termo espúrio (erro de 0,25 a 0,34\,\%).
  \item \emph{Por que o Jacobiano não é simétrico?} Pelo sinal das linhas de massa, por $S_\pi \neq \chi_\pi$
    ($\chi_w = S_w - \Gamma p_c$) e pelos termos advectivos. No FS só entra $\bm J_{uu}$, que é simétrico.
  \item \emph{Por que FS$_{GIM} \geq$ FS$_{SRM}$ no talude seco?} O Mohr--Coulomb é homogêneo de grau 1 em
    $(\bm\sigma, c)$. Aumentar o peso equivale a reduzir só a coesão; o SRM reduz $c$ e $\tan\phi$.
  \item \emph{Por que a coluna parcialmente saturada drena mais devagar?} O armazenamento da retenção,
    $nC_s \approx 2{,}7\times10^{-3}$\,kPa$^{-1}$, domina $1/M$ e $n/K_w$.
  \item \emph{O que $p_a = 0$ não testa?} O transporte do ar.
  \item \emph{Por que o SRM com $\alpha = 0$ precisava de lista?} No trecho elástico, $\bar{\bm f} = \bm 0$. O termo
    $\alpha_e^2\Delta\Lambda^2$ dá $\Delta F = l/\alpha_e$.
  \item \emph{Por que $\Delta\Lambda < $ tol pode parar cedo?} Porque $\Delta\Lambda$ depende de $l$: um passo cortado
    parece o patamar. A correção é medir no $l$ pedido.
  \item \emph{Por que FS $\geq 1$ no modelo acoplado?} Cada estado do transiente já está em equilíbrio com a resistência
    real, e todo equilíbrio é cota inferior. Um 0,992 é falha da busca.
  \item \emph{Por que $\Gamma \approx 2{,}5\,$FS antes da chuva e $\approx$ FS no fim?} Com a sucção fixa, mais peso
    aumenta também o atrito, e o GIM só vence a coesão e a sucção. Na ruptura, os dois indicam o mesmo evento.
  \item \emph{Por que FS(24\,h) tem mínimo em $k \approx I$?} Com $k < I$, a superfície satura e a infiltração fica
    limitada por $k$. Com $k > I$, a água entra toda, mas drena rápido.
  \item \emph{Por que $P_f \approx 1\,\%$ na chuva, contra 24,8\,\% no FORM com variáveis?} Com campos, a superfície de
    ruptura passa por zonas fracas e fortes e faz média. Sem campos, não há média.
  \item \emph{Por que o campo da forma fechada não depende da malha, e o que isso permite?} Porque ele é avaliado ponto
    a ponto. Os mesmos $\xi_k$ dão o mesmo campo em qualquer malha, e por isso o refinamento por realização é
    consistente.
  \item \emph{Os mesmos $\xi_k$ dão a mesma realização?} Só na mesma base: Galerkin e Nyström, sim; a forma fechada do
    retângulo, não, mas a estatística é a mesma.
  \item \emph{Por que com $P_f \approx 11\,\%$ basta Monte Carlo e com 1\,\% não?} Porque
    $\mathrm{CoV} = \sqrt{(1 - P_f)/(NP_f)}$: com 1000 realizações, 9\,\% e 31\,\%.
  \item \emph{Por que o FORM não dá a solução?} Ele troca a superfície pelo plano tangente. O erro depende da curvatura
    e não é conhecido; a amostragem por importância o corrige.
  \item \emph{Por que a amostragem por importância funciona com 6000 variáveis?} O peso depende só da projeção
    $\bm v\cdot\bm u^*$. Mas o deslocamento tem de estar certo: ruído em direções que não importam explode a variância.
\end{enumerate}

\appendix
\section{A evolução da formulação (pasta \texttt{depricated})}
\label{sec:historia}

Em ordem de data. Os zips da pasta são cópias das pastas extraídas.
\begin{enumerate}
  \item \textbf{16/10}, \path{formulacao_cam_clay}: Biot saturado e $\bm u$--$p_w$ com ar passivo, $\chi = S_r$,
    van Genuchten e Cam-Clay modificado com endurecimento por sucção.
  \item \textbf{16/10}, \path{formulacao_mohr_coulomb}: troca para Mohr--Coulomb clássico, com o retorno em valores
    principais e o SRM de Griffiths--Lane. Traz a hierarquia de modelos M0 a M4 (M2 com ar passivo, M4 Richards). Uma
    variante acrescenta o sistema adimensional, ramo que não continuou.
  \item \textbf{19/10}, \path{formulacao_mohr_coulomb_atualizado}: acrescenta a percolação permanente, verificada
    na barragem de Lam (1984) e Fredlund--Rahardjo (1993), com código em Python e Mathematica.
  \item \textbf{22/10}, \path{formulacao_usat}: recomeço só com as equações de campo de água e ar, esqueleto elástico,
    e o exemplo da barragem em regime transiente.
  \item \textbf{25/10}, versão \texttt{bousineq}: a redução 1D de Boussinesq, com ponto de saída móvel (inequação
    variacional e conjunto ativo).
  \item \textbf{26/10}, \path{biot_nao_saturado_agua_ar3}: sai Boussinesq e entra o MEF (Galerkin, Euler
    implícito, Newton). É a fonte [N] do documento atual.
  \item \textbf{28/10}, \texttt{ar4} e \path{secoes_rahardjo}: a implementação no NeoPZ (PV2, Fredlund--Xing,
    $\chi = 1$) e a reprodução do Rahardjo: 144 casos, FS inicial de 2,119, SRM por bisseção, chuva na projeção
    horizontal. Inclui os contraexemplos A1 a A6.
  \item \textbf{30/10}, \path{formulacao_fase1}: a primeira versão do documento atual, com 3 decisões em aberto. No
    mesmo dia, a versão atual aplica as decisões e acrescenta os resíduos, o Jacobiano e a verificação. Em 1/11 vem a
    versão \texttt{verick}.
\end{enumerate}
Em uma frase: de Cam-Clay com $\bm u$--$p_w$, passando pelo Mohr--Coulomb clássico com a barragem e pelo recomeço
poroelástico com água e ar, até o NeoPZ com PV2, a reprodução do Rahardjo e a formulação $\bm u$--$p_w$--$p_a$
confrontada com o código e verificada.

\section{Folha de fórmulas}
\label{sec:formulas}

\paragraph{Solo não saturado.}
\begin{gather*}
  \bm\sigma = \bm\sigma' - \alpha p_B\bm I, \quad p_B = \chi p_w + (1 - \chi)p_a, \quad p_c = p_a - p_w, \quad
  \chi_w = S_w - \Gamma p_c, \ \chi_a = S_a + \Gamma p_c \\
  S_e = [1 + (a p_c)^b]^{-m} \ \text{(van Genuchten)}, \qquad
  \Theta = [\ln(e + (\psi/a)^n)]^{-m}, \ k_w = k_s\Theta^p \ \text{(Fredlund--Xing)} \\
  \operatorname{div}\bm\sigma + \rho_b\bm g = \bm 0, \qquad \bm q_\pi = -\lambda_\pi\bm K(\nabla p_\pi - \gamma_\pi\bm g),
  \qquad \dot\zeta_\pi + \operatorname{div}\bm q_\pi + c_\pi\bm q_\pi\cdot\nabla p_\pi = 0 \\
  \Delta\zeta_w = S_w\Theta + n\Delta S_w + \tfrac{nS_w}{K_w}\Delta p_w, \quad
  \Delta\zeta_a = S_a\Theta - n\Delta S_w + \tfrac{nS_a^n}{P}\Delta p_a, \quad \Theta = \alpha\Delta\varepsilon_v +
  \eta\Delta p_B \\
  \tau_f = c' + (\sigma_n - u_a)\tan\phi' + (u_a - u_w)\tan\phi^b \quad (\chi = 1 \Leftrightarrow \phi^b = \phi')
\end{gather*}

\paragraph{Fator de segurança.}
\begin{gather*}
  \text{SRM: } c'/F,\ \tan\phi'/F,\ \tan\psi/F; \qquad
  \text{GIM: } \lambda\,(\text{peso} + \langle p_B\rangle_+), \text{ sucção fixa}; \qquad
  \text{FS} = \max\{\Lambda : \text{há equilíbrio}\} \\
  \bm K_T\delta\bm u - \bar{\bm f}\delta\Lambda = -\bm r, \qquad
  \Delta\bm u\cdot\Delta\bm u + \alpha_e^2\Delta\Lambda^2 = l^2, \qquad
  \text{parada: } \Delta\Lambda\, l_{pedido}/l < 10^{-3}
\end{gather*}

\paragraph{Chuva.} Em cada nó da superfície: $q_{in} = I$ com $p_w \leq 0$, ou $p_w = 0$ com $q_{in} \leq I$.

\paragraph{Campos aleatórios.}
\begin{gather*}
  C = e^{-|\Delta x|/L_x - |\Delta y|/L_y}, \qquad c' = e^{\mu_{\ln} + \sigma_{\ln}Z}, \quad
  \sigma_{\ln}^2 = \ln(1 + \mathrm{CoV}^2), \quad \mu_{\ln} = \ln\mu - \tfrac12\sigma_{\ln}^2 \\
  Z = \sum_{k \leq M}\sqrt{\lambda_k}\,\xi_k f_k, \qquad \int_D C f_k = \lambda_k f_k, \qquad
  \varepsilon_M = 1 - \sum_{k \leq M}\lambda_k/|D| \\
  \text{1D: } \omega\tan(\omega a) = c \ \text{(pares)}, \ \omega + c\tan(\omega a) = 0 \ \text{(ímpares)}, \
  \lambda = \tfrac{2c}{\omega^2 + c^2} \\
  \text{Nyström: } W^{1/2}KW^{1/2}\bm v = \lambda\bm v, \ f_k(\bm x_j) = v_j/\sqrt{w_j}
\end{gather*}

\paragraph{Probabilidade de falha.}
\begin{gather*}
  \hat P_f = n_f/N, \quad \mathrm{CoV} = \sqrt{(1 - P_f)/(NP_f)} \\
  \text{FORM: } \beta = \|\bm u^*\|, \ P_f \approx \Phi(-\beta), \quad
  \bm u_{k+1} = \tfrac{\nabla G_k\cdot\bm u_k - G_k}{\|\nabla G_k\|^2}\nabla G_k \\
  \text{Importância: } w = e^{-\bm v\cdot\bm u^* + \|\bm u^*\|^2/2}, \quad \hat P_f = \tfrac1N\sum I[G \leq 0]\,w \\
  \text{Nataf (LN--LN): } \rho_Z = \tfrac{\ln(1 + \rho\,\delta_i\delta_j)}{\zeta_i\zeta_j}
\end{gather*}


\end{document}
